\documentclass[11pt]{article}
\usepackage[T1]{fontenc}
\usepackage[utf8]{inputenc}
\usepackage{lmodern}
\usepackage{amsmath,amssymb,amsthm,mathtools}
\usepackage[margin=1in]{geometry}
\usepackage[hidelinks]{hyperref}
\usepackage{enumitem}

% id: proofs.integer-distance.gap-note.001
% claim_id: GAP-IDIST-CLUSTERS-01
% collection: proofs
% language: en
% source_path: proofs/claims/integer_distance_preprint.tex
% source_locator: Section "Finite integral point sets in general position"; Assumption "Unbounded integral clusters"
% proof_status: gap-present
% gap_ids: GAP-IDIST-CLUSTERS-01, GAP-IDIST-ELLIPSE-B, GAP-IDIST-HYPERBOLA-D, GAP-IDIST-PARABOLA-E

\newtheorem{theorem}{Theorem}
\newtheorem{lemma}[theorem]{Lemma}
\newtheorem{proposition}[theorem]{Proposition}
\newtheorem{corollary}[theorem]{Corollary}
\theoremstyle{definition}
\newtheorem{definition}[theorem]{Definition}
\newtheorem{question}[theorem]{Question}
\theoremstyle{remark}
\newtheorem{remark}[theorem]{Remark}

\newcommand{\C}{\mathbb{C}}
\newcommand{\Q}{\mathbb{Q}}
\newcommand{\R}{\mathbb{R}}
\newcommand{\Z}{\mathbb{Z}}
\newcommand{\T}{\mathbb{T}}
\newcommand{\ii}{\mathrm{i}}
\newcommand{\Qs}{\Q^{\square}}
\newcommand{\abs}[1]{\left|#1\right|}
\newcommand{\set}[1]{\left\{#1\right\}}

\setlist[itemize]{leftmargin=2em}
\setlist[enumerate]{leftmargin=2em}
\setlist[description]{leftmargin=2em,style=nextline}

\title{Gap Note for the Finite Integer-Distance Cluster Problem}
\author{}
\date{}

\begin{document}
\maketitle

\begin{abstract}
This proof-attempt isolates two arithmetic routes toward the finite problem left
open in \path{proofs/claims/integer_distance_preprint.tex}.  The first uses a
Joukowski map from the unit circle to a noncircular ellipse; the second uses the
rectangular hyperbola \(xy=1\).  Each route proves the geometric reduction and
leaves one explicit arithmetic gap.
\end{abstract}

\section{Provenance and status}

This note does not prove the main finite existence problem.  Its source-backed
input is \path{proofs/claims/integer_distance_preprint.tex}, section ``Finite
integral point sets in general position'', Assumption ``Unbounded integral
clusters'', also logged as \texttt{GAP-IDIST-CLUSTERS-01}.  The reductions below
are local derived claims: the geometric parts are proved here, while the terminal
arithmetic inputs remain open.

\begin{description}
\item[source-backed]
The preprint's Assumption ``Unbounded integral clusters'' asks for, for every
\(n\), an \(n\)-point planar set with pairwise integer distances, no three
collinear points, and no four concyclic points.

\item[derived-claim]
Either Missing Lemma B or Missing Lemma D below implies the source-backed
finite cluster assumption.

\item[open gap]
The arithmetic statements \texttt{GAP-IDIST-ELLIPSE-B} and
\texttt{GAP-IDIST-HYPERBOLA-D} are not proved here.  Section~\ref{sec:next-targets}
also records \texttt{GAP-IDIST-PARABOLA-E}, a new open parabola route suggested
after this note was first drafted.
\end{description}

Write
\[
\Qs=\set{q^2:q\in\Q}
\]
for the set of rational squares.

\section{The finite target}

\begin{definition}[Good finite set]
A finite set \(S\subset\R^2\) is \emph{good} if:
\begin{enumerate}
\item for any distinct \(P,Q\in S\), the distance \(\abs{P-Q}\) is a positive
integer;
\item no three points of \(S\) are collinear;
\item no four points of \(S\) are concyclic.
\end{enumerate}
\end{definition}

\begin{lemma}[Scaling rational distances]
Let \(S\subset\R^2\) be finite.  Suppose every distance between two distinct
points of \(S\) is a positive rational number.  Then there is a rational number
\(M>0\) such that the dilation \(MS=\set{MP:P\in S}\) has all pairwise distances
in \(\Z_{>0}\).  Collinearity and concyclicity relations are unchanged by this
dilation.
\end{lemma}

\begin{proof}
If \(D\), the finite set of pairwise distances in \(S\), is empty, take
\(M=1\).  Otherwise write each \(d\in D\) in lowest terms and let \(M\) be a
common multiple of all denominators.  Then \(Md\in\Z_{>0}\) for every
\(d\in D\).  A nonzero homothety maps lines to lines and circles to circles.
\end{proof}

\begin{question}[Main finite problem]
For every positive integer \(n\), does there exist a good set with exactly
\(n\) points?
\end{question}

The next two sections reduce this question to arithmetic existence problems.

\section{The Joukowski ellipse route}

Let
\[
\T=\set{z\in\C:\abs z=1}.
\]
For \(\mu\in\C^\times\), define
\[
J_\mu(z)=z+\frac{\mu}{z}\qquad (z\in\T).
\]

\begin{lemma}[Derived claim: Joukowski distance identity]
Let \(z_i,z_j\in\T\) be distinct.  Then
\[
\abs{J_\mu(z_i)-J_\mu(z_j)}
=
\abs{z_i-z_j}\,\abs{z_i z_j-\mu}.
\]
\end{lemma}

\begin{proof}
We compute
\[
J_\mu(z_i)-J_\mu(z_j)
=z_i-z_j+\mu\left(\frac1{z_i}-\frac1{z_j}\right)
=(z_i-z_j)\left(1-\frac{\mu}{z_i z_j}\right).
\]
Since \(\abs{z_i z_j}=1\), taking absolute values gives the identity.
\end{proof}

\begin{corollary}[Derived claim: rational-distance transfer]
Let \(z_1,\ldots,z_n\in\T\) be distinct and let \(\mu\in\C^\times\).  If
\[
\abs{z_i-z_j}\in\Q_{>0}
\quad\text{and}\quad
\abs{z_i z_j-\mu}\in\Q_{>0}
\]
for all \(i\ne j\), then the points \(J_\mu(z_1),\ldots,J_\mu(z_n)\) have
positive rational pairwise distances.
\end{corollary}

\begin{proof}
The distance identity gives rational positive distances, hence distinct image
points.
\end{proof}

\subsection{Ellipse normalization}

Write
\[
\mu=\rho\eta^2,
\qquad
\rho=\abs\mu>0,
\qquad
\eta\in\T.
\]
If \(z=\eta w\), then
\[
J_\mu(z)=\eta\left(w+\frac{\rho}{w}\right).
\]
Thus it is enough to study
\[
E_\rho(w)=w+\frac{\rho}{w}\qquad (w\in\T).
\]
Writing \(w=e^{\ii t}\),
\[
E_\rho(w)=(1+\rho)\cos t+\ii(1-\rho)\sin t.
\]
If \(\rho\ne 1\), this parametrizes a nondegenerate noncircular ellipse, up to
the harmless rotation by \(\eta\).

\begin{lemma}[Derived claim: no three collinear on the ellipse]
Assume \(\rho>0\) and \(\rho\ne 1\).  If \(w_1,
\ldots,w_n\in\T\) are distinct, then no three of the points
\(E_\rho(w_1),\ldots,E_\rho(w_n)\) are collinear.
\end{lemma}

\begin{proof}
The image is contained in the real conic
\[
\frac{x^2}{{(1+\rho)}^2}+\frac{y^2}{{(1-\rho)}^2}=1,
\]
whose two denominators are nonzero.  A line intersects this nondegenerate conic
in at most two real points.  Hence three distinct image points cannot be
collinear.
\end{proof}

\begin{lemma}[Derived claim: four-cyclicity on the ellipse]\label{lem:ellipse-cyclic}
Assume \(\rho>0\) and \(\rho\ne 1\).  Let
\(w_1,w_2,w_3,w_4\in\T\) be distinct.  Then the four points
\[
E_\rho(w_1),\ E_\rho(w_2),\ E_\rho(w_3),\ E_\rho(w_4)
\]
are concyclic if and only if
\[
w_1w_2w_3w_4=1.
\]
\end{lemma}

\begin{proof}
Write a general circle in complex coordinates as
\[
\zeta\overline{\zeta}+A\zeta+\overline A\,\overline\zeta+B=0,
\]
where \(A\in\C\) and \(B\in\R\).  For \(\zeta=E_\rho(w)=w+\rho w^{-1}\) and
\(\abs w=1\), we have
\[
\overline\zeta=w^{-1}+\rho w.
\]
After substitution and multiplication by \(w^2\), the circle equation becomes
\[
\rho w^4+(A+\rho\overline A)w^3+(1+\rho^2+B)w^2
+(\rho A+\overline A)w+\rho=0.
\]
Thus if four distinct parameters lie on one circle, they are the four roots of
a quartic whose leading and constant coefficients are both \(\rho\).  Their
product is therefore \(1\).

Conversely, assume \(w_1w_2w_3w_4=1\).  Let
\[
\prod_{k=1}^4(w-w_k)=w^4-e_1w^3+e_2w^2-e_3w+1.
\]
Because all \(w_k\) lie on \(\T\), conjugating is the same as inverting.  Thus
\(e_3=\overline{e_1}\).  Also \(\overline{e_2}\) is the sum of
\(\frac{1}{w_i w_j}\) over the six unordered pairs.  Since
\(w_1w_2w_3w_4=1\), these are exactly the complementary pair-products, so
\(e_2\in\R\).  Write \(e_1=a+\ii b\).  The equations
\[
A+\rho\overline A=-\rho e_1,
\qquad
\rho A+\overline A=-\rho\overline{e_1}
\]
are solved by \(A=x+\ii y\) with
\[
x=-\frac{\rho a}{1+\rho},
\qquad
y=-\frac{\rho b}{1-\rho},
\]
using \(\rho>0\) and \(\rho\ne 1\).  Set \(B=\rho e_2-1-\rho^2\), which is real.
The corresponding circle equation pulls back to
\(\rho\prod_{k=1}^4(w-w_k)=0\).  Hence the four image points lie on that circle.
\end{proof}

\begin{corollary}[Audit correction: phase in the original variables]\label{cor:phase-correction}
With the normalization \(\mu=\rho\eta^2\) and \(z=\eta w\), four points
\(J_\mu(z_a),J_\mu(z_b),J_\mu(z_c),J_\mu(z_d)\) are concyclic if and only if
\[
z_a z_b z_c z_d=\eta^4={\left(\frac{\mu}{\abs\mu}\right)}^2.
\]
\end{corollary}

\begin{proof}
By Lemma~\ref{lem:ellipse-cyclic}, concyclicity is equivalent to
\((z_a/\eta)(z_b/\eta)(z_c/\eta)(z_d/\eta)=1\).  Since
\(\mu/\abs\mu=\eta^2\), this is exactly the stated condition.
\end{proof}

\begin{theorem}[Derived claim: conditional ellipse construction]\label{thm:ellipse-conditional}
Suppose that for every \(n\ge 1\) there are distinct points
\(z_1,\ldots,z_n\in\T\) and a number \(\mu\in\C^\times\) such that:
\begin{enumerate}
\item \(\abs{z_i-z_j}\in\Q_{>0}\) for all \(i\ne j\);
\item \(\abs{z_i z_j-\mu}\in\Q_{>0}\) for all \(i\ne j\);
\item \(\abs\mu\ne 1\);
\item for all distinct \(a,b,c,d\),
\[
z_a z_b z_c z_d\ne {\left(\frac{\mu}{\abs\mu}\right)}^2.
\]
\end{enumerate}
Then good sets exist in every finite cardinality.
\end{theorem}

\begin{proof}
For a fixed \(n\), set
\[
S_n=\set{J_\mu(z_i):1\le i\le n}.
\]
The distance identity gives rational pairwise distances.  The ellipse
normalization with \(\rho=\abs\mu\ne 1\) gives no three collinear points.  The
phase criterion in Corollary~\ref{cor:phase-correction} gives no four concyclic
points.  The scaling lemma turns \(S_n\) into a good set.
\end{proof}

\section{A concrete powers specialization}

Let \(\theta\) satisfy \(\sin\theta,\cos\theta\in\Q\), and let
\[
u=e^{2\ii\theta}.
\]
Then every power of \(u\) has rational real and imaginary parts.  More is true:
the chord lengths among powers of \(u\) are rational.

\begin{lemma}[Derived claim: rational chord lengths from a rational angle point]
If \(\sin\theta,\cos\theta\in\Q\), then
\[
\abs{u^i-u^j}=2\abs{\sin((i-j)\theta)}\in\Q
\]
for all integers \(i,j\), where \(u=e^{2\ii\theta}\).
\end{lemma}

\begin{proof}
The trigonometric identity gives
\[
\abs{e^{2\ii i\theta}-e^{2\ii j\theta}}
=2\abs{\sin((i-j)\theta)}.
\]
The number \(\sin(k\theta)\) is an integral polynomial expression in
\(\sin\theta\) and \(\cos\theta\), so it is rational.
\end{proof}

\begin{remark}[Audit correction: the sample angle]
The usable Pythagorean example is
\[
\cos\theta=\frac45,
\qquad
\sin\theta=\frac35.
\]
It gives
\[
u=e^{2\ii\theta}=\frac7{25}+\frac{24}{25}\ii.
\]
The values \(5/4\) and \(5/3\) cannot be sine or cosine values of a real angle.
\end{remark}

\begin{lemma}[Derived claim: the sample power is not torsion]
For \(u=7/25+(24/25)\ii\), the unit complex number \(u\) is not a root of unity.
\end{lemma}

\begin{proof}
If a root of unity \(\zeta\) has rational real part, then
\(\zeta+\zeta^{-1}=2\operatorname{Re}(\zeta)\) is both rational and an algebraic
integer, hence an integer.  For the present \(u\), this number is
\(14/25\), not an integer.  Therefore \(u\) is not a root of unity.
\end{proof}

\begin{question}[Open gap: corrected Missing Lemma B]\label{q:missing-b}
For every \(n\ge 1\), do there exist an angle \(\theta\) with
\(\sin\theta,\cos\theta\in\Q\), a non-torsion point \(u=e^{2\ii\theta}\), and a
rational number \(\rho\in\Q_{>0}\setminus\set{1}\) such that
\[
\abs{\rho-u^m}\in\Q
\qquad
\text{for every }m=3,4,\ldots,2n-1?
\]
\end{question}

\begin{proposition}[Derived claim: Missing Lemma B implies the finite target]
If Question~\ref{q:missing-b} has a positive answer for every \(n\), then good
sets exist in every finite cardinality.
\end{proposition}

\begin{proof}
Fix \(n\), take data from Question~\ref{q:missing-b}, and set \(z_i=u^i\) for
\(1\le i\le n\).  The non-torsion condition makes the \(z_i\) distinct.  The
preceding chord-length lemma gives \(\abs{z_i-z_j}\in\Q_{>0}\).  Since
\(3\le i+j\le 2n-1\), Question~\ref{q:missing-b} gives
\(\abs{z_i z_j-\rho}\in\Q_{>0}\).  Apply Theorem~\ref{thm:ellipse-conditional}
with \(\mu=\rho\).

For the four-cyclicity obstruction, \(\mu/\abs\mu=1\).  If
\(z_a z_b z_c z_d=1\), then \(u^{a+b+c+d}=1\), contradicting non-torsion because
\(a+b+c+d>0\).  Thus the product obstruction is avoided.
\end{proof}

\begin{question}[Stronger fixed-angle target]
For the fixed value
\[
u=\frac7{25}+\frac{24}{25}\ii,
\]
does every \(n\) admit \(\rho\in\Q_{>0}\setminus\set{1}\) such that
\[
\rho^2+1-2\rho\operatorname{Re}(u^m)
\]
is a rational square for every \(m=3,4,\ldots,2n-1\)?
\end{question}

This fixed-angle target is stronger than Question~\ref{q:missing-b}.  A failure
of the fixed-angle target would not by itself refute the moving-\(u\) version.

\begin{remark}[Audit correction: rational coordinates alone]
The condition \(u\in\T\cap\Q(i)\) does not by itself imply
\(\abs{u^i-u^j}\in\Q\) for all \(i,j\).  The powers specialization above uses the
stronger hypothesis \(u=e^{2\ii\theta}\) with \(\sin\theta,\cos\theta\in\Q\), which
directly proves rational chord lengths.
\end{remark}

\section{The hyperbola route}

For \(t\in\Q^\times\), define
\[
H(t)=\left(t,\frac1t\right)\in\R^2.
\]

\begin{lemma}[Derived claim: hyperbola distance formula]
For distinct \(s,t\in\Q^\times\),
\[
\abs{H(s)-H(t)}
=
\frac{\abs{s-t}}{\abs{st}}\sqrt{{(st)}^2+1}.
\]
\end{lemma}

\begin{proof}
We have
\[
\frac1s-\frac1t=-\frac{s-t}{st}.
\]
Therefore
\[
\abs{H(s)-H(t)}^2
={(s-t)}^2+{\left(\frac1s-\frac1t\right)}^2
={(s-t)}^2\left(1+\frac1{s^2t^2}\right),
\]
which is the square of the displayed formula.
\end{proof}

\begin{corollary}[Derived claim: rational distances on the hyperbola]
Let \(t_1,
\ldots,t_n\in\Q^\times\) be distinct.  If
\[
{(t_i t_j)}^2+1\in\Qs
\]
for every \(i\ne j\), then the points \(H(t_1),\ldots,H(t_n)\) have rational
pairwise distances.
\end{corollary}

\begin{proof}
The prefactor \(\abs{t_i-t_j}/\abs{t_i t_j}\) is rational and positive.  The
remaining square-root term is rational by hypothesis.
\end{proof}

\begin{lemma}[Derived claim: no three collinear on \(xy=1\)]
No line in \(\R^2\) contains three distinct points of the hyperbola \(xy=1\).
\end{lemma}

\begin{proof}
A vertical line intersects \(xy=1\) in at most one point.  Any nonvertical line
has equation \(y=mx+b\).  Substitution gives
\[
x(mx+b)=1,
\]
a polynomial equation of degree at most two in \(x\).  Hence such a line has at
most two intersections with the hyperbola.
\end{proof}

\begin{lemma}[Derived claim: four-cyclicity on \(xy=1\)]\label{lem:hyperbola-cyclic}
Let \(t_a,t_b,t_c,t_d\in\Q^\times\) be distinct.  The four points
\[
H(t_a),\ H(t_b),\ H(t_c),\ H(t_d)
\]
are concyclic if and only if
\[
t_a t_b t_c t_d=1.
\]
\end{lemma}

\begin{proof}
A circle has equation
\[
x^2+y^2+Ax+By+C=0
\]
for some real \(A,B,C\).  Substituting \(y=1/x\) and multiplying by \(x^2\)
gives
\[
x^4+Ax^3+Cx^2+Bx+1=0.
\]
Thus, if four distinct hyperbola parameters lie on one circle, they are the
four roots of a monic quartic with constant coefficient \(1\), so their product
is \(1\).

Conversely, suppose \(t_a t_b t_c t_d=1\).  Write
\[
\prod_{r\in\set{a,b,c,d}}(x-t_r)=x^4-e_1x^3+e_2x^2-e_3x+1.
\]
Taking \(A=-e_1\), \(C=e_2\), and \(B=-e_3\), the circle equation above pulls
back to this quartic.  Hence the four corresponding points lie on a circle.
\end{proof}

\begin{corollary}[Derived claim: positivity prevents concyclicity]
If \(t_1,
\ldots,t_n\in\Q\) are all greater than \(1\), then no four of the points
\(H(t_1),\ldots,H(t_n)\) are concyclic.
\end{corollary}

\begin{proof}
The product of any four distinct selected parameters is greater than \(1\), so
Lemma~\ref{lem:hyperbola-cyclic} excludes concyclicity.
\end{proof}

\begin{question}[Open gap: Missing Lemma D]\label{q:missing-d}
For every \(n\ge 1\), do there exist distinct rationals
\[
1<t_1<\cdots<t_n
\]
such that
\[
{(t_i t_j)}^2+1\in\Qs
\qquad
\text{for all }i<j?
\]
\end{question}

\begin{theorem}[Derived claim: Missing Lemma D implies the finite target]
If Question~\ref{q:missing-d} has a positive answer for every \(n\), then good
sets exist in every finite cardinality.
\end{theorem}

\begin{proof}
For a fixed \(n\), choose \(t_1,
\ldots,t_n\) as in Question~\ref{q:missing-d} and set
\[
S_n=\set{H(t_i):1\le i\le n}.
\]
The hyperbola rational-distance corollary gives rational pairwise distances.
The hyperbola line-intersection lemma gives no three collinear points.  Since
all \(t_i>1\), the positivity corollary gives no four concyclic points.  The
scaling lemma gives a good set with \(n\) points.
\end{proof}

\section{Equivalent forms of Missing Lemma D}

Question~\ref{q:missing-d} can be restated in several ways useful for later
agents.

\begin{proposition}[Derived claim: square-entry form]
Question~\ref{q:missing-d} is equivalent to the following statement: for every
\(n\ge 1\), there exist distinct rational squares
\[
a_1,
\ldots,a_n>1
\]
such that
\[
a_i a_j+1\in\Qs
\qquad
\text{for all }i<j.
\]
\end{proposition}

\begin{proof}
If \(t_i\) solve Question~\ref{q:missing-d}, set \(a_i=t_i^2\).  Conversely, if
\(a_i\) are rational squares greater than \(1\), choose their positive rational
square roots \(t_i>1\).  The displayed conditions are then identical.
\end{proof}

\begin{proposition}[Derived claim: multiplicative clique form]
Let
\[
P=\set{x\in\Q_{>0}:x^2+1\in\Qs}.
\]
Question~\ref{q:missing-d} is equivalent to asking for arbitrarily large finite
sets \(T\subset\Q_{>1}\) such that
\[
st\in P
\qquad
\text{for all distinct }s,t\in T.
\]
Moreover,
\[
P=\set{\frac{1-q^2}{2q}:q\in\Q,
0<q<1}
=\set{\frac{2r}{r^2-1}:r\in\Q,
r>1}.
\]
\end{proposition}

\begin{proof}
The equivalence with Question~\ref{q:missing-d} follows from taking
\(T=\set{t_1,
\ldots,t_n}\).

For the parametrization, solve \(x^2+1=y^2\) with \(x,y\in\Q\) and \(x>0\).
Then
\[
(y-x)(y+x)=1.
\]
Put \(q=y-x\).  Since \(x>0\) and \(y>0\), one has \(0<q<1\), and
\[
x=\frac{1-q^2}{2q}.
\]
Conversely, this formula gives a positive rational \(x\) with
\(x^2+1={((1+q^2)/(2q))}^2\).  The second parametrization follows by the change
of variable \(q=(r-1)/(r+1)\), where \(r>1\).
\end{proof}

\section{Gap register}

\begin{description}
\item[\texttt{GAP-IDIST-CLUSTERS-01}]
Source-backed finite obstruction from
\path{proofs/claims/integer_distance_preprint.tex}: construct, for every
\(n\), an \(n\)-point planar integral cluster with no three collinear points and
no four concyclic points.

\item[\texttt{GAP-IDIST-ELLIPSE-B}]
Open arithmetic input for the ellipse route: prove Question~\ref{q:missing-b},
or prove a variant strong enough to satisfy Theorem~\ref{thm:ellipse-conditional}.
This is a simultaneous rational-distance problem from one rational real point
to finitely many powers on the unit circle.

\item[\texttt{GAP-IDIST-HYPERBOLA-D}]
Open arithmetic input for the hyperbola route: prove Question~\ref{q:missing-d}.
Equivalently, construct arbitrarily large rational Diophantine tuples whose
entries are rational squares.

\item[\texttt{GAP-IDIST-PARABOLA-E}]
Open arithmetic input from the parabola route in Section~\ref{sec:next-targets}:
construct arbitrarily large finite sets of positive rationals whose pairwise
sums are rational legs of rational right triangles.

\item[\texttt{GAP-IDIST-JOUKOWSKI-PHASE-01}]
Audit item resolved in Corollary~\ref{cor:phase-correction}: under the
normalization \(\mu=\rho\eta^2\), the product obstruction in the original
\(z\)-variables is \({(\mu/\abs\mu)}^2\).

\item[\texttt{GAP-IDIST-UNIT-POWERS-01}]
Audit item resolved in the powers specialization: rational real and imaginary
parts of \(u\) alone are not used as the chord-length hypothesis.  The note uses
the stronger rational-angle condition \(u=e^{2\ii\theta}\) with
\(\sin\theta,\cos\theta\in\Q\).
\end{description}

\section{Next attack target: standalone problem}\label{sec:next-targets}

The most promising first gap to attack is now \texttt{GAP-IDIST-PARABOLA-E}.
It is not proved here, and the proposed arithmetic-progression/CRT argument for
it should not be treated as source-backed.  Its advantage is that the geometric
verification is elementary and completely separated from the arithmetic input.

\begin{question}[Open gap: standalone parabola problem]\label{q:missing-e}
For every integer \(n\ge 1\), do there exist distinct positive rational numbers
\[
x_1,\ldots,x_n
\]
such that
\[
1+{(x_i+x_j)}^2\in\Qs
\qquad
\text{for all }i<j?
\]
Equivalently, can one find arbitrarily large finite subsets of \(\Q_{>0}\)
whose pairwise sums are rational numbers \(s\) for which \(1+s^2\) is a rational
square?
\end{question}

\begin{proposition}[Derived claim: the parabola problem implies the finite target]
If Question~\ref{q:missing-e} has a positive answer for every \(n\), then good
sets exist in every finite cardinality.
\end{proposition}

\begin{proof}
Given \(x_1,\ldots,x_n\) from Question~\ref{q:missing-e}, set
\[
P_i=(x_i,x_i^2).
\]
For \(i\ne j\),
\[
\abs{P_i-P_j}^2
={(x_i-x_j)}^2\bigl(1+{(x_i+x_j)}^2\bigr),
\]
so the pairwise distances are rational.  A line meets the parabola \(y=x^2\) in
at most two points, so no three selected points are collinear.  If four selected
points lie on a circle
\[
x^2+y^2+Dx+Ey+F=0,
\]
then substituting \(y=x^2\) gives
\[
x^4+(E+1)x^2+Dx+F=0.
\]
The four corresponding positive \(x\)-coordinates would be the roots of this
quartic, whose \(x^3\)-coefficient is zero; their sum would be zero, impossible.
The scaling lemma then gives a good set.
\end{proof}

A useful but stronger subproblem is to find, for every \(n\), a nonzero rational
number \(d\) such that
\[
1+m^2d^2\in\Qs
\qquad (m=3,4,\ldots,2n-1).
\]
Then \(x_i=id\) would solve Question~\ref{q:missing-e}.  This stronger statement
is currently only a candidate attack route; the quoted Chinese-remainder or
Kemnitz-style construction is not verified in this note and is false as stated
for the sample formula tested during audit.

The hyperbola target \texttt{GAP-IDIST-HYPERBOLA-D} remains the next cleanest
alternative: find rational numbers \(t_i>1\) with
\[
{(t_i t_j)}^2+1\in\Qs
\]
for all pairs.  The ellipse target \texttt{GAP-IDIST-ELLIPSE-B} offers more
parameters, but its bookkeeping is less convenient for a first standalone
attack.

\end{document}
